\chapter{Código-fonte do infográfico}
\label{anx:codigo}

\section{Repositório completo}

O código-fonte completo do infográfico, com histórico de \emph{commits}
e \texttt{README} de instruções, encontra-se em:

\begin{center}
\fbox{\parbox{0.9\textwidth}{\centering\vspace{6pt}
\large \url{https://github.com/glkwolff/Projeto-Estruturas-de-computadores}
\vspace{6pt}}}
\end{center}

A reprodução integral neste anexo é impraticável: o arquivo
\texttt{projeto.html} ultrapassa 2\,500 linhas e o critério de
entrega prevê justamente que, em código longo, se descreva o link
no relatório. As listagens a seguir trazem apenas trechos críticos
ilustrativos.

\section{Estrutura do arquivo único}

\begin{lstlisting}[language=HTML,caption=Esqueleto geral do projeto.html]
<!DOCTYPE html>
<html lang="pt-BR">
<head>
  <meta charset="UTF-8">
  <title>Arquitetura de Servidores - Infográfico Interativo</title>
  <style>
    /* ~600 linhas de CSS:
       custom properties, layout flex/grid, tabelas, animações */
  </style>
</head>
<body>
  <header> ... </header>
  <div class="container">
    <nav>  <!-- 7 botões de navegação --> </nav>
    <main>
      <section id="mod1"  class="module active"> <!-- Módulo 1 --> </section>
      <section id="mod2"  class="module">        <!-- Módulo 2 --> </section>
      <section id="mod3"  class="module">        <!-- Módulo 3 --> </section>
      <section id="mod4"  class="module">        <!-- Módulo 4 --> </section>
      <section id="mod5"  class="module">        <!-- Módulo 5 --> </section>
      <section id="mod6"  class="module">        <!-- Módulo 6 --> </section>
      <section id="modRef" class="module">       <!-- Referências --> </section>
    </main>
  </div>
  <footer> ... </footer>
  <script>
    /* ~600 linhas de JavaScript:
       navegação, sub-abas, simulações por módulo */
  </script>
</body>
</html>
\end{lstlisting}

\section{Função de navegação por sub-abas}

\begin{lstlisting}[language=JavaScript,caption=Implementação das sub-abas dentro de cada módulo]
function showPanel(modId, panelId, element) {
    const root = document.getElementById(modId);
    if (!root) return;
    root.querySelectorAll('.mod-panel').forEach(p =>
        p.classList.remove('active'));
    root.querySelectorAll('.mod-tab').forEach(t =>
        t.classList.remove('active'));
    const panel = document.getElementById(panelId);
    if (panel) panel.classList.add('active');
    if (element) element.classList.add('active');
}
\end{lstlisting}

\section{Cálculo do simulador RAID}

\begin{lstlisting}[language=JavaScript,caption=Calculadora de IOPS e capacidade útil]
function calcRaid() {
    const level = parseInt(document.getElementById('rc-level').value);
    let   n     = parseInt(document.getElementById('rc-disks').value);
    const cap   = parseFloat(document.getElementById('rc-cap').value);
    const iops  = parseInt(document.getElementById('rc-iops').value);

    const mins = { 0: 2, 1: 2, 5: 3, 6: 4, 10: 4 };
    if (n < mins[level]) n = mins[level];
    if (level === 10 && n % 2 !== 0) n -= 1;

    let useful, tolerable, penalty;
    switch (level) {
        case 0:  useful = n * cap;       tolerable = 0;     penalty = 1; break;
        case 1:  useful = cap;           tolerable = n-1;   penalty = 2; break;
        case 5:  useful = (n-1) * cap;   tolerable = 1;     penalty = 4; break;
        case 6:  useful = (n-2) * cap;   tolerable = 2;     penalty = 6; break;
        case 10: useful = (n/2) * cap;   tolerable = n/2;   penalty = 2; break;
    }
    const iopsR = n * iops;
    const iopsW = Math.round((n * iops) / penalty);
    /* ... atualização da DOM omitida ... */
}
\end{lstlisting}

\section{Aplicação da Lei de Amdahl}

\begin{lstlisting}[language=JavaScript,caption=Simulador interativo da Lei de Amdahl]
function updateAmdahl() {
    const P = parseFloat(document.getElementById('am-p').value);
    const N = parseInt  (document.getElementById('am-n').value);

    const speedup = 1 / ((1 - P) + (P / N));
    const ceiling = P === 1 ? Infinity : 1 / (1 - P);
    const util    = ceiling === Infinity ? '—'
                  : ((speedup / ceiling) * 100).toFixed(0) + '%';

    document.getElementById('am-out-speedup').innerText = speedup.toFixed(2) + 'x';
    document.getElementById('am-out-cap').innerText     = ceiling === Infinity
                                                        ? 'inf'
                                                        : ceiling.toFixed(2) + 'x';
    document.getElementById('am-out-util').innerText    = util;
}
\end{lstlisting}

\section{Simulador de \textit{cache hit/miss}}

\begin{lstlisting}[language=JavaScript,caption=Animação de travessia na hierarquia de memória]
function simulateCacheAccess() {
    const r = Math.random();
    let hitAt, cycles;
    if      (r < 0.92)  { hitAt = 'l1';  cycles = 4;   }
    else if (r < 0.98)  { hitAt = 'l2';  cycles = 12;  }
    else if (r < 0.995) { hitAt = 'l3';  cycles = 38;  }
    else                { hitAt = 'ram'; cycles = 250; }

    const order = ['l1', 'l2', 'l3', 'ram'];
    const hitIndex = order.indexOf(hitAt);
    order.forEach((lvl, i) => {
        setTimeout(() => {
            const el = document.querySelector(
                `#cache-trace .cache-level[data-cl="${lvl}"]`);
            if (i <  hitIndex) el.classList.add('miss');
            if (i === hitIndex) el.classList.add('hit');
        }, i * 250);
    });
}
\end{lstlisting}

\section{Navegação para referência citada}

\begin{lstlisting}[language=JavaScript,caption=Função que vincula cita inline ao item bibliográfico]
function goToRef(n) {
    showModule('modRef', null);
    setTimeout(() => {
        const target = document.getElementById('ref-' + n);
        if (!target) return;
        target.scrollIntoView({ behavior: 'smooth', block: 'center' });
        target.classList.add('highlight');
        setTimeout(() => target.classList.remove('highlight'), 2500);
    }, 200);
}
\end{lstlisting}

\section{Como executar localmente}

\begin{enumerate}
    \item Clonar o repositório: \\
          \texttt{git clone https://github.com/glkwolff/Projeto-Estruturas-de-computadores}
    \item Abrir o arquivo \texttt{projeto.html} em qualquer navegador
          moderno (Chrome 130+, Firefox 130+, Edge 130+).
    \item Nenhuma instalação, servidor local ou conexão de internet é
          necessária.
\end{enumerate}
